% Immutable cited-input supplement. Target proofs remain embargoed.
% Target: inv:wb:uniform-nontorsion-provenance; tag 01TQ.
% Frozen source SHA256 74497bfe327fae25c9cdcb17334b3bf42a9793a9df7c0a4404d97568f2f760f0.
% CompendiumRef is resolved against the cited S6_09 source edition, not silently identified with the current series statement.

% Exact cited input: series:I / inv:A1.0
\begin{cited}[Structure-sheaf direct images]\hypertarget{tag-01RQ}{}\label{inv:A1.0}\label{wb:thm:structure-direct-images-audit}\label{wb:lem:cusp-structure-sheaf-audit}


This is the structure-sheaf calculation of Claude and Alp\"oge: \cite[Lemma~9.12, pp.~72--74; Proposition~9.13, pp.~74--76, and the base-change clarification on p.~92]{AC2026}. The direct images $R^qf_*\OO_X$, for $q=0,1,2$, are
$\OO$, $G=\OO\gamma\oplus\OO(-1)$, and $\OO(-1)$.
The inclusion $\OO\gamma\subset G$ is a subbundle.  At the cusp,
$(\gamma,u)$ is a holomorphic frame, and cup product identifies
$\bigwedge^2G$ with $R^2f_*\OO_X$.  On the normalized root family, the cohomology pullback by the
cyclic generator has characters $(1,\zeta^{-a})$, where
$(m,a)=(3,2)$ or $(4,1)$ and $\zeta=e^{-2\pi i/m}$.
These assertions concern full scheme fibres, including $3S_1$ and $4S_2$.
At the cusp, the natural map $H^1(W;\mathbb C)\to H^1(W,\mathcal O_W)$ is an isomorphism. In degree two the comparison is surjective and is injective on $\ker[H^2(W;\mathbb C)\to H^2(\mathrm{dP}_6;\mathbb C)]$. The cup product $\bigwedge^2H^1(W,\mathcal O_W)\to H^2(W,\mathcal O_W)$ is an isomorphism. Every base-change map $(R^qf_*\mathcal O_X)\otimes k(b)\to H^q(X_b,\mathcal O_{X_b})$ is an isomorphism, including at the nonreduced fibres.
\end{cited}

% Exact cited input: comp:I / mapsI:top:thm:p0-integral-specialization
\begin{theorem}[nearby-cycle specialization]
\label{mapsI:top:thm:p0-integral-specialization}
\ For the proper holomorphic cusp filling
\(f_0:N_0\to\Delta\) of \cref{wb:thm:candidate-cusp-geometry}, with its
proved local normal-crossings equations and normalization combinatorics,
the specialization map of the central fibre \(W\),
\[
 \operatorname{sp}^r:H^r(W;\mathbb Z)\longrightarrow H^r(F;\mathbb Z)
\]
is injective and
\[
 \operatorname{im}\operatorname{sp}^r
 =H^r(F;\mathbb Z)^{T_0}
\tag{IL.5.1}
\]
for \(0\le r\le4\).  In particular, the image is saturated and there is no
unrecorded finite-index defect at \(p_0\).
\end{theorem}

% Exact cited input: comp:I / wb:thm:candidate-cusp-geometry
\begin{theorem}[the cusp quotient and its central fibre]
\label{wb:thm:candidate-cusp-geometry}
\ For the period functions and holomorphic cusp matrix \(C(t)\)
constructed above, and sufficiently small \(\epsilon>0\), the maps
\[
 \Psi_{\bar\lambda}(x,t)=
 \bigl(c_{\bar\lambda}(t)t^{B_0\bar\lambda}x,t\bigr),
 \qquad
 c_{\bar\lambda}(t)=\exp(2\pi iC(t)\bar\lambda),
\tag{2.5}
\]
extend to a free properly discontinuous action of
\(\bar\Lambda\simeq\mathbb Z^2\) on
\(Y_\epsilon=\{|t|<\epsilon\}\).  The quotient
\[
 f_0:N_0=Y_\epsilon/\bar\Lambda\longrightarrow\Delta_\epsilon
\]
is a proper holomorphic map from a connected smooth complex threefold.
Its nonzero fibres are the tori of \cref{wb:prop:period-input-audit}.  Its full
scheme-theoretic central fibre \(W\) is reduced and has local equations
\[
 z_0,\qquad z_0z_1,\qquad z_0z_1z_2.
\tag{2.6}
\]
The surface \(W\) is irreducible and nonnormal.  Its normalization is
\(\nu:E_0\simeq\mathrm{dP}_6\to W\), obtained by identifying opposite sides
of the toric hexagon; its double locus consists of three rational curves
through exactly two triple points.  Thus \(W\) is a normal-crossings surface,
but it is not a simple-normal-crossings divisor.  Moreover
\[
 \omega_{N_0}\simeq\mathcal O_{N_0}.
\tag{2.6a}
\]
\end{theorem}

% Exact cited input: comp:I / wb:prop:period-input-audit
\begin{proposition}[existence and non-degeneracy of the period matrix]
\label{wb:prop:period-input-audit}
\ There are holomorphic functions
\[
 \tau:\mathfrak h_z\to\mathfrak h,
 \qquad \mu,\beta:\mathfrak h_z\to\mathbb C
\]
satisfying
\[
\begin{array}{lll}
\tau\circ g_0=\tau-1,&\mu\circ g_0=\mu,&\beta\circ g_0=\beta+1,\\
\tau\circ g_2=-1/\tau,&\mu\circ g_2=1+\mu/\tau,
 &\beta\circ g_2=\beta-3-6\mu^2/\tau,
\end{array}
\tag{2.1}
\]
and the corresponding order-three transformation laws obtained from
\(\tau\circ g_1=(\tau-1)/\tau\).  They may be chosen so that
\[
 \operatorname{Im}\tau>0,
 \qquad
 D:=\operatorname{Im}\beta-
 \frac{6(\operatorname{Im}\mu)^2}{\operatorname{Im}\tau}<0.
\tag{2.2}
\]
Consequently
\[
 \Pi=[Z\mid I_2],\qquad
 Z=\begin{pmatrix}6\mu&\tau\\ \beta&\mu\end{pmatrix}
\tag{2.3}
\]
maps the marked rank-four lattice \(\Lambda\) isomorphically onto a lattice
in \(\mathbb C^2\) at every point.  The quotients
\(\mathbb C^2/\Pi(z)\Lambda\) form a holomorphic family of compact
two-dimensional complex tori over \(\mathfrak h_z\), equivariant for
\(\Delta\).
\end{proposition}

